\documentclass[11pt]{article}
\usepackage[margin=25mm]{geometry}
\usepackage{fontspec}
\setmainfont{Latin Modern Roman}
\usepackage{amsmath,amssymb,mathtools}
\usepackage{unicode-math}
\setmathfont{Latin Modern Math}
\usepackage{xurl,hyperref}
\hypersetup{colorlinks=true,urlcolor=blue}
\setcounter{secnumdepth}{0}
\setlength{\parindent}{0pt}
\setlength{\parskip}{5pt}
\setlength{\emergencystretch}{3em}
\newcommand{\sref}[2]{\hyperlink{source:#1:#2}{[S#2]}}
\newcommand{\cataloguescope}[1]{\paragraph{Recorded target.}#1}
\begin{document}
\section{The Odd Perfect Number Conjecture}\label{statement}

\subsection{Definitions and mathematical statement}

\paragraph{Shared notation.}$\mathbb N=\{1,2,\ldots\}$, and $\mathbb Z,\mathbb Q,\mathbb R,\mathbb C$ denote the integers, rationals, reals and complex numbers. Any other conventions are specified locally.


For \(n\in\mathbb N\), let
\(\sigma(n)=\sum_{d\mid n,\ d>0}d\) be the sum of all its positive
divisors.  The integer \(n\) is perfect when \(\sigma(n)=2n\),
equivalently when it equals the sum of its proper positive divisors.
\paragraph{Statement recorded in the supplied source.}
Is there any odd integer \(n>1\) satisfying
\[
 \sigma(n)=2n?
\]
Equivalently, is the set of odd perfect numbers empty? \sref{statement}{2}
\subsection{Short English statement}
Does an odd perfect number exist?
\subsection{Sources}
\begin{description}
\item[\hypertarget{source:statement:1}{S1}] ULAM.AI and the UnsolvedMath Contributors, \emph{UnsolvedMath: A Curated Collection of Open Mathematics Problems}, catalogue entry. CC BY 4.0. \url{https://huggingface.co/datasets/ulamai/UnsolvedMath}.
\item[\hypertarget{source:statement:2}{S2}] Pascal Ochem and Joshua Zelinsky, \emph{On odd perfect numbers with exactly one even exponent greater than 2}, arXiv:2607.19746 (2026), Section 1, definition of an odd perfect number. \url{https://arxiv.org/abs/2607.19746}.
\end{description}
\label{end:statement}
\end{document}
